\chapter{写在开始之前}\label{写在开始之前}

这本教材邀请我们共同完成一件事：把一个善于续写文本的大语言模型（Large Language Model，LLM），逐步变成能讲题、能查资料、能使用工具，也能承认信息不足的学习助手。我们暂时叫它``小助''。名字只是为了叙述方便。

我们会保留一条清楚的故事线。起初，小助连``请回答问题''都不总能做到。我们用案例示范教它回答，它才开始像一个助手，却会在续训新课程资料后丢掉原来的讲解习惯。我们让它多试几次，发现偶尔能答对，于是尝试筛选和强化学习。分数开始上涨，我们又发现验证集没有进步。最后，我们把检索、工具和评估接入系统，才意识到：可靠的助手来自模型、数据、训练与应用共同作用。

我们默认自己会写基本程序，知道数组和函数的含义。正文首次引入术语时给出中文名、英文全称与常用缩写，后文沿用；模型名和论文题名保留原名或给出代称。遇到概率、向量、梯度时，正文会就地补上够用的数学。初读时，我们先追踪``问题为什么出现、证据支持什么、下一步如何验证''；第二遍再推公式和复现实验。全书按九部分展开，适合约八到十二周的学习，也可以按课堂节奏拆分。每个编号小节都包含完整讲解，以及关键论文与补充材料的阅读入口。

书中``小助''的回答、行为和数值案例均为教学构造，用来练习推理与诊断。在实际实验或者工程任务中，则需要记录模型版本、数据划分、模板、随机种子、计算预算与解码配置。

本材料一到三章希望构建一个对大模型的基本理解，包括tokenizer，词嵌入，句子嵌入，并且把大模型的行为理解为 next token prediction的条件分布。四到六章串讲现代大模型全部训练方法。后续章节则尽量对大模型各个方向做一个入门性质的学习。

任何方法都不保证在所有模型和任务上有效；我们关心的是如何识别和理解任务中的有效条件。本教材尽可能不堆砌概念和知识点逐一介绍，而是尽可能拉通每个部分，希望读者能够自然地理解从上一步到下一步的逻辑关系，这样有助于理解整个大模型体系。本教材省略了许多技术细节，这些技术细节是成为大模型专家不可或缺的一部分，如果读者立志成为该领域的专家，建议读者自学 Hugging Face 大语言模型课程与 Stanford CS336《从零构建语言模型》的视频和材料。

实际上，在当前的科研甚至工业研发场景下，本教材所提到的许多知识可能永远都用不上，例如一名后训练算法工程师通常只需要掌握数据合成、后训练方法（最基础的SFT和RL）就可以完成大部分工作，甚至不需要关心词表、嵌入层等模型架构的知识，一名大模型工程师真正重要的能力在于碰到问题时能够有思路去摸索和解决问题，这就要求我们有尽可能准确和发散的思维，有科学可复用的方法论指导实践。